\cotefile{46}
\scanpage{24}

\begin{center}
Definition of the Verschiebung morphism.
\end{center}

I do not know exactly what you want me to define\deleted{\unsure{z|s}}. However, since you claim to have lost my carpet from last year, I am going to try to sell it to you again. For this, I am going to suppose that the base scheme is the spectrum of a field $k$ of characteristic $p\handwritten{>}0$, perfect or not: it is up to you to see what happens in the other cases.

I am going to begin with the case of a commutative \uline{affin\deleted{\unsure{t|é}}e} $k$-group $G$, with coordinate ring $\mathbb{A}$. The notation is as follows:

\noindent - $\mathbb{A}^{(p)}$ is the $k$-algebra obtained from $\mathbb{A}$ by extension of scalars $F_k:k\mathrel{\text{-{}-{}-{}-{}-}\!\handwritten{>}}k$, where $\mathbb{F}_k(x)=x^p$. An element of $\mathbb{A}^{(p)}$ will therefore be written $x\handwritten{\otimes}a$, where $x\handwritten{\in}k$ and $a\handwritten{\in}\mathbb{A}$; we shall note the calculation rule $x\handwritten{\otimes}y.a=x.y^p\handwritten{\otimes}a$. If $k$ is perfect, $\mathbb{A}^{(p)}$ is identified with the $k$-algebra obtained from $\mathbb{A}$ by restriction of scalars $x\mathrel{\text{-{}-{}-{}-{}-}\!\handwritten{>}}x^{1\handwritten{/}p}$.

\noindent - $F:\mathbb{A}^{(p)}\mathrel{\text{-{}-{}-{}-{}-}\!\handwritten{>}}\mathbb{A}$ is the map \deleted{\illegible{1}} $x\handwritten{\otimes}a\mathrel{\text{-{}-{}-{}-{}-}\!\handwritten{>}}x.a^p$.

\noindent - $G^{(p)}$ is the algebraic group with coordinate ring $A^{(p)}$; we shall also denote by $F$ the morphism from $G$ to $G^{(p)}$ which is defined by the map $F:\mathbb{A}^{(p)}\mathrel{\text{-{}-{}-{}-{}-}\!\handwritten{>}}A$.

\noindent - $\handwritten{\otimes}^{p}A$ denotes the tensor product of $p$ copies of $A$; $\handwritten{\Delta}^{p}$ is the map from $\mathbb{A}$ to $\handwritten{\otimes}^{p}\mathbb{A}$ which is defined by the diagonal $\handwritten{\Delta}$ of the coalgebra $\mathbb{A}$; $m^p$ is
% archive sheet number: 23
